\chapter{よくあるエラーと対処}
\label{chap:errors}

\section{\cmd{uv: command not found}}

\textbf{原因}：uv がインストールされていないか、PATH が通っていません。

\textbf{対処1}：ターミナルを閉じて開き直す。

\textbf{対処2}：PATH を今すぐ反映させる。

\begin{wslbox}
\begin{lstlisting}[style=bash]
$ source ~/.bashrc
\end{lstlisting}
\end{wslbox}

\textbf{対処3}：それでも解決しない場合は PATH を手動で追加する。

\begin{wslbox}[PATH を手動で追加]
\begin{lstlisting}[style=bash]
$ echo 'export PATH="$HOME/.local/bin:$PATH"' \
  >> ~/.bashrc
$ source ~/.bashrc
\end{lstlisting}
\end{wslbox}

最後に確認：

\begin{wslbox}[確認]
\begin{lstlisting}[style=bash]
$ uv --version
uv 0.6.x (linux x86_64 ...)
\end{lstlisting}
\end{wslbox}

\section{\cmd{No pyproject.toml found}}

\textbf{原因}：\cmd{pyproject.toml} のあるフォルダにいません。

\textbf{対処}：正しいフォルダに移動してから実行する。

\begin{wslbox}[正しいフォルダへ移動して実行]
\begin{lstlisting}[style=bash]
$ cd ~/NMR-lab-kadai/2026_B4_assignment
$ uv sync
\end{lstlisting}
\end{wslbox}

フォルダに \cmd{pyproject.toml} があるか確認：

\begin{wslbox}[pyproject.toml の確認]
\begin{lstlisting}[style=bash]
$ ls pyproject.toml
pyproject.toml
\end{lstlisting}
\end{wslbox}

\cmd{pyproject.toml} と表示されれば正しいフォルダにいます。

\section{JupyterLab でカーネルが見つからない}

\textbf{原因}：\cmd{uv run} をつけずに JupyterLab を起動しています。

\textbf{対処}：ターミナルで \key{Ctrl} + \key{C} を押して一度終了し、
正しいコマンドで起動し直す。

\begin{wslbox}[正しい起動コマンド]
\begin{lstlisting}[style=bash]
$ cd ~/NMR-lab-kadai/2026_B4_assignment
$ uv run jupyter lab
\end{lstlisting}
\end{wslbox}

\section{\cmd{ModuleNotFoundError} が出る}

\textbf{原因1}：\cmd{uv sync} が完了していない。

\textbf{対処}：\cmd{uv sync} を再実行する。

\begin{wslbox}
\begin{lstlisting}[style=bash]
$ cd ~/NMR-lab-kadai/2026_B4_assignment
$ uv sync
\end{lstlisting}
\end{wslbox}

\textbf{原因2}：\cmd{uv run} なしで JupyterLab を起動している。
$\rightarrow$ 前節の対処と同じ手順で起動し直す。

\section{\cmd{uv sync} がネットワークエラーで止まる}

\textbf{原因}：インターネット接続が不安定、または
学内ネットワークのプロキシが干渉しています。

\textbf{対処}：数分待ってから再実行する。
\cmd{uv sync} は途中まで進んだ場合、次回実行時に続きから再開します。

\begin{wslbox}[再実行]
\begin{lstlisting}[style=bash]
$ uv sync
\end{lstlisting}
\end{wslbox}

繰り返し発生する場合は教員・先輩に相談してください。

\section{VS Code で「カーネルを起動できませんでした」と表示される}

\textbf{症状}：ノートブックを開くと
「Python 環境 'nmrlab-b4-2026 (3.12.x)' を使用できなくなったため、
カーネルを起動できませんでした」と表示される。

\textbf{原因}：VS Code が以前選択したカーネルをキャッシュしているため、
\cmd{.venv/} を作り直したあとに古い参照が残ります。
\cmd{pyvenv.cfg} を手動で書き換えても \cmd{uv sync} で元に戻るので、
この方法では解決しません。

\textbf{対処}：\cmd{uv sync} で \cmd{.venv/} を作り直したあと、
VS Code でカーネルを選び直す。

\begin{enumerate}
  \item ターミナルで \cmd{.venv/} を削除して再構築する
\end{enumerate}

\begin{wslbox}[.venv を削除して再構築]
\begin{lstlisting}[style=bash]
$ cd ~/NMR-lab-kadai/2026_B4_assignment
$ rm -rf .venv
$ uv sync
\end{lstlisting}
\end{wslbox}

\begin{enumerate}
  \setcounter{enumi}{1}
  \item VS Code でノートブックを開き、右上の
    「カーネルを選択」をクリック
  \item 「別のカーネルを選択」$\rightarrow$「Python 環境」を選ぶ
  \item 一覧から \cmd{.venv/bin/python}（または
    \texttt{nmrlab-b4-2026 (3.12.x)}）を選択する
\end{enumerate}

\begin{notebox}[一覧に出ない場合]
「Python 環境の一覧を更新」をクリックするか、
VS Code を開き直してから再度選択してください。
\end{notebox}

\section{Jupyter で「Python 3.12.x が使用できなくなった」と表示される}

\textbf{原因}：フォルダ内に古い \cmd{.venv/} が残っており、
そこに保存されていた Python のパスが現在の環境に存在しません。
（配布ファイルを古いバージョンから上書き展開した場合に起こります。）

\textbf{対処}：\cmd{.venv/} を削除してから \cmd{uv sync} を再実行する。

\begin{wslbox}[.venv を削除して再構築]
\begin{lstlisting}[style=bash]
$ cd ~/NMR-lab-kadai/2026_B4_assignment
$ rm -rf .venv
$ uv sync
\end{lstlisting}
\end{wslbox}

\cmd{uv sync} が完了したら VS Code でフォルダを開き直し、
カーネルを \cmd{.venv/bin/python} に選び直してください。

\section{\cmd{uv sync} 後に \cmd{Permission denied} が出る}

\textbf{原因}：フォルダを Windows の共有ドライブから直接コピーした場合、
ファイルの実行権限が失われることがあります。

\textbf{対処1}：\cmd{.venv/} を削除して再構築する（最も確実）。

\begin{wslbox}[.venv を削除して再構築]
\begin{lstlisting}[style=bash]
$ cd ~/NMR-lab-kadai/2026_B4_assignment
$ rm -rf .venv
$ uv sync
\end{lstlisting}
\end{wslbox}

\textbf{対処2}：フォルダ全体の権限を修正する。

\begin{wslbox}[権限の変更]
\begin{lstlisting}[style=bash]
$ chmod -R u+rw \
  ~/NMR-lab-kadai/2026_B4_assignment
\end{lstlisting}
\end{wslbox}

\begin{notebox}[git clone で取得した場合は発生しにくい]
\cmd{git clone} で正しく取得した場合、パーミッション問題は
ほとんど起きません。問題が続く場合は \cmd{.venv/} を削除して
\cmd{uv sync} を再実行してください（第6章参照）。
\end{notebox}

\section{ブラウザが自動で開かない}

\textbf{原因}：WSL から Windows のブラウザを自動起動できない場合があります。

\textbf{対処}：ターミナルに表示された URL を手動でコピーしてブラウザに貼り付ける。

\begin{wslbox}[URL の場所]
\begin{lstlisting}[style=bash]
[I ServerApp] Jupyter Server is running at:
[I ServerApp] http://localhost:8888/lab
  ?token=abc123def456...  <- これを全部コピーする
\end{lstlisting}
\end{wslbox}
